\section{小规模受控实测：核验门禁的缺陷检出率}\label{sec:experiment}

\subsection{为什么要做，以及做到哪一步}

第 \ref{sec:metrics} 节把反例检出率 $\mathrm{CDR}$ 定义为“注入的带缺陷样本中被验证流程拦截的比例”。该定义若只停在定义式，就只是一个\emph{可计算性}的声明，无法说明一个具体验证流程究竟拦得住什么、拦不住什么。本节做一次最小规模的实例化：在一份\emph{冻结}的稿件快照上做受控缺陷注入，实测一个真实门禁的检出率、基线误报与盲区。

需要先说清口径：$\mathrm{CDR}$ 的原始定义针对\emph{规约缺陷}（漏写假设、扩大定义域、偷偷加入正性假设）。本节实测的对象是\emph{引文核验门禁}，因此这是同一方法论（注入已知缺陷、测量拦截比例）在一个完全机械可判定对象上的应用，而不是对 $\mathrm{CDR}$ 本身的测量；对规约缺陷的同类测量需要先把四类判据机械化，留作后续工作（第 \ref{sec:mutex} 小节）。选择引文核验作为首个实测对象有三个理由：它是十项清单中判据完全可机械化的部分；它对应学术写作中的一条明确红线（编造或错配引用）；它的全部输入都可以离线冻结，因而实验确定、可复跑、不需要调用任何模型。

\subsection{设置与被检物}

被检物由三条互相独立的信息源构成：（1）正文引用集合——\texttt{parts/*.tex} 中 \texttt{\textbackslash cite} 命令展开后的 key 集合；（2）条目库——\texttt{references.bib}，共 31 条；（3）登记快照——\texttt{lit/refs.json}，记录 31 条条目在 arXiv 与 Crossref 的查询结果（标题、年份、来源 URL），实验期间只读冻结，实验过程不联网。

门禁脚本 \texttt{pipeline/check\_refs\_static.py} 在上述三个输入上做确定性比对：一组结构规则（悬空引用、幽灵条目、重复 key、DOI 格式、登记存在性、标题相似度、年份一致性），配标题相似度阈值 0.85（FAIL）与 0.95（WARN）、年份容差 $\pm 1$ 年。退出码构成契约：\texttt{0} 表示通过，\texttt{20} 表示存在 FAIL 级发现。

\paragraph{基线。}对\emph{未改动}的稿件裸跑门禁，返回 \textbf{0 条 FAIL 级发现}（输出 \texttt{RESULT: PASS}）。这一条是后面所有数字的解释前提：既然基线没有任何既有告警，那么后续每一处命中都只能归因于注入本身，而“误报率 0”正是指基线零发现。

\paragraph{注入方式。}\texttt{pipeline/defect\_injection.py} 按类别逐条构造变体，\emph{每次只改动一处}，然后把变体交给同一门禁函数，检查被判定的 key 是否恰好是目标条目。八个臂对应八类可机械表达的缺陷；每条真实条目在每类下贡献一个注入样本。

\begin{table}[t]
\centering
\small
\caption{引文核验门禁的受控缺陷注入结果（单一冻结快照，31 条条目）。D6 不注入代码级变体：语义错配无法由机械规则表达，故该臂按“每条目一个潜在样例”计数，拦截数由设计决定为 0。}
\label{tab:injection}
\begin{tabularx}{\textwidth}{@{}X c c c c@{}}
\toprule
\textbf{缺陷类别} & \textbf{期望规则} & \textbf{注入} & \textbf{拦截} & \textbf{检出率} \\
\midrule
D1 悬空引用（引用不存在的条目） & R1 & 31 & 31 & 100\% \\
D2 幽灵条目（库中有、正文未引） & R2 & 31 & 31 & 100\% \\
D3a DOI 格式非法 & R5 & 12 & 12 & 100\% \\
D3b 编造 DOI（格式合法、登记不存在） & R6 & 12 & 6 & 50\% \\
D4a 标题张冠李戴（与相邻条目轮换） & R7 & 31 & 31 & 100\% \\
D4b 年份漂移 $+7$ & R7 & 31 & 31 & 100\% \\
D5 重复 key（追加同名条目） & R3 & 31 & 31 & 100\% \\
D6 语义错配（条目真实、主张错配） & 无 & (31) & 0 & 0\% \\
\midrule
\textbf{结构类合计（D1--D5）} & --- & \textbf{210} & \textbf{173} & \textbf{82.4\%} \\
\bottomrule
\end{tabularx}
\end{table}

\subsection{结果}

八类缺陷共注入 210 个样本，拦截 173 个，\textbf{结构类检出率 82.4\%}；基线误报 0 条。分歧只出现在两个臂上。

\paragraph{D3b 编造 DOI：50\%。}十二个样本中有六个漏检。这一臂的构造是“把 DOI 换成一个格式合法、但登记快照中查不到的字符串”，因此它的检出依赖门禁能否对“这个 DOI 是否真实存在”给出判断。本文只报告整体数字，未逐条归因；一个待验证的假设是：漏检集中在快照中缺少 DOI 对照记录的条目上，此时门禁无对照可比，只能放过。若该假设成立，那么这一分歧不是判据缺陷，而是\textbf{离线快照设计的固有代价}——不联网就无法区分“查不到”与“不存在”。这恰好给出一条工程建议：对 DOI 这类存在性断言，核验阶段应允许一次联网解析来换取这一类覆盖率。

\paragraph{D6 语义错配：0\%，且这是设计上的 0。}“条目真实、标题与年份都对，只是被挂在错误的主张上”是一种结构完全合法的缺陷：门禁看到的每条引用都指向真实存在的文献，因此没有任何一条结构规则可以触发。本文把这个臂计入分母，而不是把它排除在外——排除它会让表观检出率变成 100\%，恰好掩盖最需要被看见的盲区。这一臂的 0 说明一件事：\emph{机械核验能保证“引用的对象真实”，不能保证“引用支撑了该主张”}。后者只能由人工逐条做“标题—主张字面比对”来拦截。在本稿的定稿过程中，这一人工环节实际抓出 4 处引用—主张错配（详见随附的 AI 流程日志），若只依赖门禁它们将全部进入送审稿。这正是本文“机械化到哪一层、哪一层必须留给人”这一区分的具体形状。

\subsection{该实测能支持什么，不能支持什么}

三点限制必须与数字一起读。第一，\textbf{构造偏差}：注入缺陷由脚本按规则生成，其分布天然偏向门禁可识别的形态；这一偏差无法通过增加样本量消除，只能靠加入检出率为 0 的盲区臂把它显式化（D6 的作用）。第二，\textbf{规模}：样本是单一稿件快照（31 条条目），分歧臂只有 12 个样本，因此这些数应被读作\emph{该门禁行为的刻画}，而不是对检出率的总体估计。第三，\textbf{对象外推}：结论只适用于这类“输入可冻结、判据可机械表达”的核验环节；把它外推到模型能力或数据集表现是超出支持的。

\paragraph{核验是流程里的硬门禁，不是事后检查。}本文把引用核验接到编译之前：\texttt{pipeline/build.py} 的第一环就是核验脚本，它非零退出时第二环（编译）不会被执行。因此“核验通过才能产出 PDF”不是自我承诺，而是可以从退出码逐环核对的执行顺序——任一引用登记失败，\texttt{paper.pdf} 就不会被更新；反过来，只要产物是新的，核验在这一次构建里就是过关的。

为使“非零即中止”不被读作自我承诺，\texttt{pipeline/gate\_negative\_test.py} 把这条执行顺序做成对照实验：同一环境、同一脚本跑两份副本，未变异副本 \texttt{result=PASS} 且报告中存在编译环；把一条未登记引用注入正文后，核验环退出码 20、整体 \texttt{result=FAIL}，报告中\emph{不存在}编译环、输出里也不出现编译环节点。门禁因此是被观测到的执行顺序，而不是事后检查（留痕：\texttt{out/experiment/gate-negative-test.txt}）。

复跑方式（四条命令，无需网络）：

\begin{quote}
\texttt{python3 pipeline/build.py} \quad 一次走完硬门禁：核验（非零即中止，不进入编译）$\rightarrow$ 编译 $\rightarrow$ 产物复核，末行打印 \texttt{RESULT: PASS/FAIL}；\\
\texttt{python3 pipeline/check\_refs\_static.py} \quad 只跑核验环，期望输出 \texttt{RESULT: PASS}；\\
\texttt{python3 pipeline/defect\_injection.py} \quad 重放八个注入臂，打印上表与合计检出率；\\
\texttt{python3 pipeline/gate\_negative\_test.py} \quad 负向验证：变异输入下编译环确实不执行，期望 \texttt{RESULT: PASS}。
\end{quote}

逐臂结果同时以结构化形式落盘在 \texttt{out/experiment/defect-injection.json}，与可读报告 \texttt{out/experiment/defect-injection.md} 并置，便于他人复核而不必重跑。
